\documentclass[11pt,a4paper]{article}

\usepackage[T1]{fontenc}
\usepackage[utf8]{inputenc}
\usepackage{amsmath,amssymb,amsthm}
\usepackage{booktabs}
\usepackage{geometry}
\geometry{margin=2.7cm}
\usepackage{xcolor}
\usepackage[colorlinks=true,linkcolor=blue!50!black,citecolor=blue!50!black]{hyperref}

\newtheorem{proposition}{Proposition}
\newtheorem{lemma}{Lemma}
\newtheorem{corollary}{Corollary}
\theoremstyle{remark}
\newtheorem{remark}{Remark}
\newtheorem{assumption}{Assumption}

\newcommand{\R}{\mathbb{R}}
\newcommand{\PP}{\mathbb{P}}
\newcommand{\dd}{\,\mathrm{d}}
\newcommand{\norm}[1]{\left\lVert #1 \right\rVert}
\newcommand{\nK}[1]{\left\lVert #1 \right\rVert_{K}}
\newcommand{\ip}[2]{\left( #1 , #2 \right)}
\newcommand{\symgrad}{\nabla^{\mathrm{S}}}
\newcommand{\Proj}{\Pi}
\newcommand{\chat}{\hat{c}}

\title{A note on the dimension and shape dependence of the viscous
stabilization constant $c_1$ for quadratic simplicial elements in
ASGS/VMS methods}
\author{Guillermo Casas\thanks{CIMNE -- Centre Internacional de M\`etodes
Num\`erics en Enginyeria, Barcelona. Draft note; comments welcome.}}
\date{\today}

\begin{document}
\maketitle

\begin{abstract}
In numerical experiments with an algebraic subgrid-scale (ASGS)
stabilized finite element method for the Brinkman/Navier--Stokes
equations, we observed that quadratic tetrahedral meshes require a
value of the algorithmic constant $c_1$ (the coefficient of
$\nu/h_K^2$ in the inverse stabilization parameter) roughly four times
larger than quadratic triangular meshes in order to obtain stable,
optimally convergent solutions, whereas the standard practice inherited
from the literature is to use the same constants in two and three
dimensions. We show that this observation is consistent with, and
quantitatively predicted by, the structure of the stability analysis:
the elementwise coercivity of the Galerkin-plus-stabilization energy
imposes the lower bound $c_1 \ge c_1^*(K) = 2\,\chat^2(K)$, where
$\chat^2(K)$ is the largest eigenvalue of an explicit local
generalized eigenvalue problem involving the divergence of the
projected velocity gradient. We compute $c_1^*(K)$ exactly for the
relevant simplex shapes. For $\PP_1$ elements $c_1^*(K)=0$ identically
--- the mechanism is absent, which explains why it has left no trace in
the (predominantly linear-element) stabilized-methods literature. For
$\PP_2$ elements the threshold is strongly shape- and
dimension-dependent; for the element shapes produced by structured
mesh generation the predicted ratio between tetrahedra and triangles is
$(100+5\sqrt{2})/24 \approx 4.46$, and across the shape pairs
considered the ratio lies in $[2.8,\,4.5]$. The ingredients
(inverse-inequality constants as design parameters) are classical, but
we have not found the two- to three-dimensional rescaling of $c_1$
stated explicitly anywhere.
\end{abstract}

\section{The observation and the question}
\label{sec:observation}

We consider the stabilized finite element approximation of the
incompressible Brinkman/Navier--Stokes equations in a variable-porosity
medium, discretized with continuous $\PP_k$ velocity--pressure pairs
and stabilized with an ASGS/VMS term-by-term method. The stabilization
parameters follow the standard design of Codina and
coworkers~\cite{Codina2000,Codina2002,CodinaVMSChapter}: on each
element $K$ of diameter $h_K$,
%
\begin{equation}
\label{eq:taus}
\tau_{1,K}
 = \Bigl( \alpha_K \tau_{\mathrm{NS},K}^{-1} + \sigma \Bigr)^{-1},
 \qquad
\tau_{\mathrm{NS},K}^{-1}
 = c_1 \frac{\nu}{h_K^2} + c_2 \frac{|\boldsymbol a|_K}{h_K},
 \qquad
\tau_{2,K} = \frac{h_K^2\,\tau_{\mathrm{NS},K}^{-1}}{c_1\,\alpha_K},
\end{equation}
%
where $\alpha_K$ is an elementwise representative of the porosity,
$\sigma$ the reaction (Darcy) coefficient, and $c_1$, $c_2$ the
algorithmic constants, for which the standard values $c_1 = 4$, $c_2 =
2$ (with $h_K$ divided by the polynomial order for higher-order
interpolations) are ubiquitous in that
literature~\cite{Codina2000,Codina2002,CodinaVMSChapter}, in two and
three dimensions alike. The viscous term of the momentum operator is
written with a pointwise orthogonal projection $\Proj$ of the velocity
gradient; in our formulation $\Proj$ extracts the
\emph{deviatoric--symmetric} part,
%
\begin{equation}
\label{eq:proj}
\Proj \nabla \boldsymbol u
 \;=\; \symgrad \boldsymbol u
   - \tfrac{1}{d}\,(\nabla\!\cdot\!\boldsymbol u)\,\mathbb I ,
\qquad
\symgrad \boldsymbol u
 = \tfrac12\bigl(\nabla\boldsymbol u + \nabla\boldsymbol u^{\!\top}\bigr),
\end{equation}
%
so that the viscous operator reads
$-2\nabla\!\cdot\!(\alpha\nu\,\Proj\nabla\boldsymbol u)$ and the
stabilization subtracts, elementwise, the $\tau$-weighted pairing of
the adjoint and direct residuals in the usual ASGS fashion.

In convergence studies with manufactured solutions we found the
following. With $\PP_2$ elements, a value of $c_1$ that yields stable
and optimally convergent results on triangular meshes fails on
tetrahedral meshes; stability and optimal rates are recovered on the
tetrahedra only after enlarging $c_1$ by a factor of approximately
four. With $\PP_1$ elements the same constants work in both dimensions.
We were not aware of any statement in the VMS/ASGS literature
prescribing a dimension-dependent $c_1$, so the questions are:
\emph{is the observation consistent with the theory, why is the factor
of the observed size, and why has it not been reported?}

The answer developed below is that all three questions are resolved by
a single, explicitly computable local quantity: the sharp constant of
the inverse estimate for the operator
$\boldsymbol u \mapsto \nabla\!\cdot(\Proj\nabla\boldsymbol u)$
on a single element, which (i) vanishes identically for $\PP_1$
velocities, and (ii) for $\PP_2$ velocities is between $2.8$ and $4.5$
times larger (in the squared, $c_1$-relevant normalization) on the
usual tetrahedron shapes than on the usual triangle shapes.

\section{The mechanism: elementwise coercivity of the viscous ledger}
\label{sec:mechanism}

Throughout this section, fix an element $K$ and consider the
\emph{diffusion-dominated, reaction-free limit} on that element:
$\boldsymbol a = \boldsymbol 0$, $\sigma = 0$, and $\alpha$ constant on
$K$. This is the binding regime for the constraint derived below (see
Remark~\ref{rem:relax}), and it is also the regime probed by
fine-mesh convergence studies at moderate Reynolds number. Since the
projector~\eqref{eq:proj} is a pointwise algebraic operation, on each
element $\Proj\nabla\boldsymbol u_h$ is computed exactly from
$\nabla\boldsymbol u_h$, with no auxiliary $L^2$ projection involved.

Testing the stabilized form with the solution itself, the contribution
of the element $K$ to the total energy is, in this limit, the exact
elemental identity
%
\begin{equation}
\label{eq:ledger}
Q_K(\boldsymbol u_h)
 \;=\;
 2\alpha\nu \nK{\Proj\nabla\boldsymbol u_h}^2
 \;-\; \tau_{1,K}\,\nK{\,2\nabla\!\cdot\!(\alpha\nu\,\Proj\nabla
        \boldsymbol u_h)\,}^2
 \;+\; \tau_{2,K}\,\nK{\nabla\!\cdot\!(\alpha\boldsymbol u_h)}^2 .
\end{equation}
%
Indeed, with $\boldsymbol a = \boldsymbol 0$, $\sigma = 0$ and $p_h =
q_h = 0$ the momentum residual and its adjoint both reduce to
$-2\nabla\!\cdot(\alpha\nu\Proj\nabla\boldsymbol u_h)$, so the
momentum stabilization contributes the (negative) middle term, while
the mass residual and its adjoint are $\pm\nabla\!\cdot(\alpha
\boldsymbol u_h)$, contributing the (positive) last term; the Galerkin
viscous term gives the first, using that $\Proj$ is a pointwise
orthogonal projection. With $\boldsymbol a = \boldsymbol 0$, $\sigma =
0$ and $\alpha$ constant, \eqref{eq:taus} gives the exact values
$\tau_{1,K} = h_K^2/(c_1\alpha\nu)$ and $\tau_{2,K} = \nu/\alpha$;
substituting, all physical parameters factor out:
%
\begin{equation}
\label{eq:ledger2}
\frac{Q_K(\boldsymbol u_h)}{\alpha\nu}
 \;=\;
 2 \nK{\Proj\nabla\boldsymbol u_h}^2
 - \frac{4 h_K^2}{c_1}\,
   \nK{\nabla\!\cdot\!(\Proj\nabla\boldsymbol u_h)}^2
 + \nK{\nabla\!\cdot\!\boldsymbol u_h}^2 .
\end{equation}
%
Since the restriction to $K$ of a conforming $\PP_k$ velocity field is
an arbitrary element of $[\PP_k(K)]^d$, the sign of the elemental
ledger is characterized exactly by a local eigenvalue problem.

\begin{proposition}[elementwise coercivity threshold]
\label{prop:threshold}
Let $k \ge 1$ and define
%
\begin{equation}
\label{eq:c1star}
c_1^*(K)
 \;\coloneqq\;
 \max_{\substack{\boldsymbol u \in [\PP_k(K)]^d \\
                  \Proj\nabla\boldsymbol u \neq 0}}
 \frac{4 h_K^2 \nK{\nabla\!\cdot\!(\Proj\nabla\boldsymbol u)}^2}
      {2\nK{\Proj\nabla\boldsymbol u}^2
       + \nK{\nabla\!\cdot\!\boldsymbol u}^2}.
\end{equation}
%
Then $Q_K(\boldsymbol u_h) \ge 0$ for all
$\boldsymbol u_h\vert_K \in [\PP_k(K)]^d$ if and only if
$c_1 \ge c_1^*(K)$; and if $c_1 < c_1^*(K)$ there exists an elemental
velocity field for which $Q_K < 0$ while
$\nK{\Proj\nabla\boldsymbol u_h} > 0$. Moreover, introducing
%
\begin{equation}
\label{eq:chat}
\chat^{\,2}_{\mathrm{df}}(K)
 \coloneqq h_K^2 \max_{\substack{\boldsymbol u \in [\PP_k(K)]^d\\
   \nabla\cdot\boldsymbol u = 0,\ \symgrad\boldsymbol u \neq 0}}
 \frac{\nK{\nabla\!\cdot\!\symgrad\boldsymbol u}^2}
      {\nK{\symgrad\boldsymbol u}^2},
 \qquad
\chat^{\,2}_{\mathrm{all}}(K)
 \coloneqq h_K^2 \max_{\substack{\boldsymbol u\\ \Proj\nabla
   \boldsymbol u\neq 0}}
 \frac{\nK{\nabla\!\cdot\!(\Proj\nabla\boldsymbol u)}^2}
      {\nK{\Proj\nabla\boldsymbol u}^2},
\end{equation}
%
one has the two-sided bound
$2\,\chat^{\,2}_{\mathrm{df}}(K) \le c_1^*(K) \le
 2\,\chat^{\,2}_{\mathrm{all}}(K)$.
All three quantities are the largest eigenvalues of generalized
eigenvalue problems with exactly assemblable Gram matrices on
$[\PP_k(K)]^d$, and they depend on $K$ only through its shape (they are
invariant under similarity transformations of $K$).
\end{proposition}

\begin{proof}
The equivalence and the strict-negativity statement are immediate
from~\eqref{eq:ledger2}: $Q_K \ge 0$ for all elemental fields iff the
Rayleigh quotient in~\eqref{eq:c1star} never exceeds $c_1$, i.e.\ iff
$c_1 \ge c_1^*(K)$; a maximizer realizes $Q_K < 0$ when $c_1 <
c_1^*(K)$. For the lower bound in the two-sided estimate, restrict the
maximization in~\eqref{eq:c1star} to pointwise divergence-free fields:
there $\nabla\!\cdot\!\boldsymbol u = 0$ kills the last denominator
term, and $\Proj\nabla\boldsymbol u = \symgrad\boldsymbol u$
by~\eqref{eq:proj}, giving the value
$2\chat^{\,2}_{\mathrm{df}}$. For the upper bound, drop the
(nonnegative) last denominator term. Both extremal problems are
generalized eigenvalue problems $A\boldsymbol c = \lambda B\boldsymbol
c$ for the Gram matrices of the linear maps $\boldsymbol u \mapsto
\nabla\!\cdot(\Proj\nabla\boldsymbol u)$ etc.\ in any basis of
$[\PP_k(K)]^d$ (restricted, in the divergence-free case, to the
kernel of the linear divergence constraints). Shape invariance: under
$x \mapsto \lambda R x + b$ with $R$ orthogonal, gradients scale by
$\lambda^{-1}$ and conjugate by $R$, second derivatives by
$\lambda^{-2}$; every quotient in \eqref{eq:c1star}--\eqref{eq:chat}
therefore scales by $\lambda^{-2}$, exactly compensated by $h_K^2$.
\end{proof}

\begin{proposition}[the mechanism is absent for linear elements]
\label{prop:P1}
For $k = 1$ on any element, in any dimension,
$\nabla\!\cdot\!(\Proj\nabla\boldsymbol u) \equiv 0$ for all
$\boldsymbol u \in [\PP_1(K)]^d$; hence
$c_1^*(K) = \chat^{\,2}_{\mathrm{df}}(K) =
\chat^{\,2}_{\mathrm{all}}(K) = 0$, and the elemental
ledger~\eqref{eq:ledger2} is nonnegative for every $c_1 > 0$.
\end{proposition}

\begin{proof}
For linear $\boldsymbol u$ the gradient is constant, hence so is
$\Proj\nabla\boldsymbol u$, and its divergence vanishes.
\end{proof}

\begin{remark}[elementwise versus global]
\label{rem:global}
Proposition~\ref{prop:threshold} is a statement about the elemental
energy ledger, i.e.\ about the quantity that the standard coercivity
analysis controls term by term. It does not by itself produce a
\emph{global} discrete velocity field with negative total energy when
$c_1 < c_1^*(K)$, because inter-element continuity constrains which
elemental worst modes can coexist; empirically, the practical
stability threshold sits somewhat below $c_1^*$. What the elementwise
analysis does predict robustly is the \emph{ratio} of thresholds
between element families, since the conformity slack, the
$h$-convention of the implementation, and the polynomial-order
convention are common to the two- and three-dimensional runs and
cancel in the ratio. This ratio is precisely what the experiment
measures when the same code, constants and conventions are used in 2D
and 3D.
\end{remark}

\begin{remark}[why the diffusive limit is binding]
\label{rem:relax}
For $|\boldsymbol a| > 0$ or $\sigma > 0$, $\tau_{1,K}$
in~\eqref{eq:taus} decreases, which weakens the negative term
in~\eqref{eq:ledger}; the convective and reactive couplings introduce
new cross terms, but in the sufficient-condition analysis these are
absorbed by a Young inequality whose cost again scales with the
\emph{same} inverse-estimate constants (see
Section~\ref{sec:consistency}). The diffusion-dominated elemental
constraint of Proposition~\ref{prop:threshold} is therefore the sharp
edge of the requirement, and it is the regime reached by mesh
refinement at fixed physical parameters.
\end{remark}

\section{Exact values of the threshold for the relevant shapes}
\label{sec:values}

By Proposition~\ref{prop:threshold} the threshold is a pure shape
constant. We computed the three quantities of
\eqref{eq:c1star}--\eqref{eq:chat} for $\PP_2$ velocities on the five
simplex shapes most relevant in practice, by exact rational assembly
of the Gram matrices and generalized eigenvalue solves; the results
were validated by (i) exact rational eigenvalue computation where the
matrices are rational (the right triangle), (ii) $50$-digit
arithmetic against closed forms elsewhere, and (iii) Monte--Carlo
Rayleigh sampling never exceeding the reported maxima. The
computation scripts accompany this note. In all five cases we found
$\chat^{\,2}_{\mathrm{df}} = \chat^{\,2}_{\mathrm{all}}$ (the
maximizer is divergence-free once the deviatoric
projector~\eqref{eq:proj} is used), so that the two-sided bound of
Proposition~\ref{prop:threshold} collapses and
%
\begin{equation}
\label{eq:collapse}
c_1^*(K) \;=\; 2\,\chat^{\,2}(K),
\qquad
\chat^{\,2}(K) \coloneqq \chat^{\,2}_{\mathrm{df}}(K)
             = \chat^{\,2}_{\mathrm{all}}(K).
\end{equation}
%
Note also that on divergence-free $\PP_2$ fields
$\nabla\!\cdot\symgrad\boldsymbol u = \tfrac12\Delta\boldsymbol u$ is a
\emph{constant} vector, which is what makes exact evaluation
straightforward.

\begin{table}[ht]
\centering
\caption{Sharp elementwise constants for $\PP_2$ velocities with the
deviatoric--symmetric projector~\eqref{eq:proj}; $h_K$ is the element
diameter. ``Right'' denotes the simplex with mutually orthogonal unit
legs (structured quadrilateral/hexahedral splitting also produces the
Kuhn shape in 3D). Exact values; decimal forms rounded.}
\label{tab:constants}
\begin{tabular}{lccc}
\toprule
Element shape $K$ & $\chat^{\,2}(K)$ & $c_1^*(K) = 2\chat^{\,2}(K)$ &
$c_1^*(K)$ (decimal) \\
\midrule
2D right triangle (quad split) & $48$ & $96$ & $96.0$ \\
2D equilateral triangle        & $24$ & $48$ & $48.0$ \\
3D right tetrahedron           & $440/3$ & $880/3$ & $293.3$ \\
3D Kuhn tetrahedron (cube split) & $200+10\sqrt2$ & $400+20\sqrt2$ &
$428.3$ \\
3D regular tetrahedron         & $200/3$ & $400/3$ & $133.3$ \\
\midrule
any shape, $\PP_1$ velocities  & $0$ & $0$ & $0$ \\
\bottomrule
\end{tabular}
\end{table}

The dimensional comparison, i.e.\ the predicted factor by which $c_1$
must be enlarged when moving from triangles to tetrahedra \emph{with
the same code, constants and conventions}, is the ratio of thresholds
for comparable element quality:
%
\begin{equation}
\label{eq:ratios}
\frac{c_1^*(\text{Kuhn tet})}{c_1^*(\text{right tri})}
 = \frac{100+5\sqrt2}{24} \approx 4.46,
\qquad
\frac{c_1^*(\text{right tet})}{c_1^*(\text{right tri})}
 = \frac{55}{18} \approx 3.06,
\qquad
\frac{c_1^*(\text{regular})}{c_1^*(\text{equilateral})}
 = \frac{25}{9} \approx 2.78 .
\end{equation}
%
Unstructured tetrahedral meshes of ordinary quality contain elements
between the regular and the Kuhn shapes (and worse), and the binding
element is the worst one, so the practically relevant window for the
3D/2D threshold ratio is roughly $[2.8,\,4.5]$, with the structured
pairing at $4.46$.

\begin{remark}[$h$-convention sensitivity]
\label{rem:hconv}
Table~\ref{tab:constants} uses $h_K = \operatorname{diam} K$, as in the
convergence analysis. If the implementation instead uses the grid
spacing $\ell$ of the background structured mesh (a common choice),
the entries rescale by $(\ell/h_K)^2$, which is $1/2$ for the right
triangle and $1/3$ for the Kuhn tetrahedron, so the structured-pairing
ratio becomes $\tfrac23 \cdot 4.46 \approx 2.97$; a mean-edge
convention gives an intermediate value $\approx 3.6$. The observed
factor of $\approx 4$ is thus consistent with the prediction under the
diameter convention, and in any case lies inside the window spanned by
the standard conventions and shape pairs. What no convention changes
is the conclusion that the same $c_1$ cannot be optimal for triangles
and tetrahedra with $\PP_2$ elements.
\end{remark}

\section{Consistency with the a priori analysis}
\label{sec:consistency}

The observation is also consistent with the sufficient-condition side
of the theory. In the stability analysis of the method (for general
$\sigma \ge 0$, $\boldsymbol a \neq \boldsymbol 0$ and variable
porosity), coercivity of the stabilized form in the natural mesh-de\-pendent
norm is established under the condition
%
\begin{equation}
\label{eq:sufficient}
c_1 \;>\; 2\,\xi\,\bar C^2, \qquad \xi > 2,
\end{equation}
%
where $\xi$ is the free parameter of the Young inequality used to
absorb the reactive--viscous coupling and $\bar C$ is the constant of
the weighted local inverse estimate
$\nK{\nabla\!\cdot(\alpha\,\Proj\nabla\boldsymbol u_h)} \le (\bar
C/h_K)\,\alpha_K^{1/2} \nK{\alpha^{1/2}\Proj\nabla\boldsymbol u_h}$.
(We have verified this chain of implications, at the level of every
algebraic step, in a machine-checked formal development accompanying
the underlying paper; the condition~\eqref{eq:sufficient} with the
explicit constants is not an artifact of loose estimation but the
exact requirement of the proof.) In particular
$c_1 > 4\bar C^2$ is \emph{necessary for the sufficient condition},
and for constant porosity the sharp value of $\bar C^2$ on an element
is exactly $\chat^{\,2}_{\mathrm{all}}(K)$ of~\eqref{eq:chat}. The
sufficient-condition threshold therefore carries the same shape
constants and the same 3D/2D ratios as the necessary elementwise
threshold of Proposition~\ref{prop:threshold}: the two sides of the
theory bracket the practical threshold and predict the same
dimensional rescaling. For variable porosity, the analysis further
decomposes $\bar C = \sqrt{d\,\delta_\alpha}\,C_{\mathrm{inv}} +
C_\alpha$ (product rule plus the componentwise bound
$\norm{\nabla\!\cdot M} \le \sqrt d\,\norm{\nabla M}$, with
$\delta_\alpha$, $C_\alpha$ the porosity-resolution constants), which
adds an \emph{explicit} factor $d$ on top of the shape dependence of
$C_{\mathrm{inv}}$; this route gives 3D/2D ratios between $2.5$ and
$4.3$ for the same shape pairs, again consistent with the window
of~\eqref{eq:ratios}.

\section{Relation to the literature}
\label{sec:literature}

The two ingredients of the mechanism are classical, but we have been
unable to find the conclusion --- that $c_1$ must be rescaled by a
factor of order $3$--$4.5$ from triangles to tetrahedra for quadratic
elements --- stated anywhere; the following seems to be an accurate
summary of the state of the literature (based on a targeted but
non-exhaustive search).

\emph{(i) Inverse-estimate constants as design parameters.} That
stabilization parameters must respect local inverse-inequality
constants, and that these constants depend on the element type,
polynomial degree and dimension, is the central message of Harari and
Hughes~\cite{HarariHughes1992}, who tabulate sharp constants for this
purpose. In the Galerkin/least-squares and SUPG analyses of Franca and
coworkers the guard appears explicitly as the factor $m_k =
\min\{1/3,\,2C_k\}$ in the stabilization parameter, with $C_k$ the
constant of the inverse estimate for the second-order term
\cite{FrancaFrey1992}; this is exactly the same mechanism as
Proposition~\ref{prop:threshold}, in a different formulation. Explicit
dimension-dependent constants for trace-type inverse inequalities on
simplices are given by Warburton and
Hesthaven~\cite{WarburtonHesthaven2003}. In the residual-based VMS
tradition of Hughes and coworkers, the parameter is built from the
element metric tensor, $\tau_M^{-2} \ni C_I\,\nu^2\, G\!:\!G$, so that
element shape, size anisotropy and dimension enter \emph{automatically}
through $G\!:\!G$, with $C_I$ an inverse-estimate constant computed per
element type~\cite{Bazilevs2007}. In that framework the present
phenomenon is priced in by construction.

\emph{(ii) The Codina-school parameter design.} The scalar,
$h_K$-based design~\eqref{eq:taus} with fixed $c_1 = 4$, $c_2 = 2$
(and $h_K/k$ for higher orders) is standard across two and three
dimensions in the orthogonal- and algebraic-subscale
literature~\cite{Codina2000,Codina2002,CodinaVMSChapter}, and we found
no statement there of a dimension-dependent adjustment.
Proposition~\ref{prop:P1} explains why none was ever needed: the
analyses and the vast majority of the computations in that line are
carried out with linear (or bilinear) elements, for which the
constraining eigenvalue vanishes identically, so \emph{any} positive
$c_1$ keeps the elemental viscous ledger nonnegative and the choice of
$c_1$ is governed by accuracy rather than by coercivity. The
constraint materializes only for $k \ge 2$, where the second-order
part of the residual is active elementwise --- and it is then strongly
dimension- and shape-dependent, as Table~\ref{tab:constants} shows.
We also note that the $h_K/k$ convention, designed around the
first-order (convective) scaling of the residual, compensates the
degree dependence of the \emph{second-order} viscous term only
partially, which is a separate reason why higher-order elements sit
closer to the coercivity edge than linear ones at the same nominal
$c_1$.

\emph{(iii) This paper's formulation.} Two features of the present
method sharpen the effect relative to the classical Navier--Stokes
ASGS: the deviatoric--symmetric projection~\eqref{eq:proj} (which
changes the constants but not, as our computations show, the
divergence-free extremal mechanism), and the porosity weighting, whose
analysis contributes the explicit additional factor $d$ noted in
Section~\ref{sec:consistency}.

\section{Practical recommendations}
\label{sec:recommendations}

\begin{enumerate}
\item \textbf{Dimension- and shape-aware $c_1$.} If a value
$c_1^{2\mathrm D}$ has been validated on triangles, use on tetrahedra
$c_1^{3\mathrm D} = \rho\, c_1^{2\mathrm D}$ with $\rho$ taken from
the appropriate row ratio of Table~\ref{tab:constants}; for typical
mesh pairs $\rho \in [2.8, 4.5]$, and $\rho \approx 4.5$ for
structured meshes under the diameter convention. This matches the
empirically found factor of $\approx 4$.
\item \textbf{Per-element calibration.} Alternatively, compute
$c_1^*(K)$ once per element type (or per element, for anisotropic
meshes) from the local generalized eigenvalue
problem~\eqref{eq:c1star}: assemble the two Gram matrices of
$\nabla\!\cdot(\Proj\nabla\cdot)$ and of the denominator on
$[\PP_k(\hat K)]^d$ and take the largest generalized eigenvalue --- a
one-off, reference-element-sized computation --- and set $c_1 = s\,
c_1^*(K)$ with a modest safety factor $s$. This is the
metric-tensor/$C_I$ philosophy of~\cite{HarariHughes1992,Bazilevs2007}
transplanted to the scalar-$\tau$ design.
\item \textbf{Degree scaling.} For $k \ge 2$ do not rely on the
$h_K/k$ substitution to protect the viscous term; the relevant
operator is second order and its constant grows faster than $k^2$.
Either of the two previous items handles the degree dependence
correctly.
\end{enumerate}

\section*{Reproducibility}
The exact values in Table~\ref{tab:constants} were obtained with a
short Python/SymPy script (exact rational Gram assembly; generalized
eigenvalue solves; validation by exact rational eigencomputation for
the rational-vertex cases, $50$-digit arithmetic against the closed
forms elsewhere, and Monte--Carlo Rayleigh sampling), which is
distributed together with this note.

\begin{thebibliography}{9}

\bibitem{HarariHughes1992}
I.~Harari and T.~J.~R.~Hughes,
\emph{What are $C$ and $h$? Inequalities for the analysis and design
of finite element methods},
Comput.\ Methods Appl.\ Mech.\ Engrg.\ 97 (1992) 157--192.

\bibitem{FrancaFrey1992}
L.~P.~Franca and S.~L.~Frey,
\emph{Stabilized finite element methods: II. The incompressible
Navier--Stokes equations},
Comput.\ Methods Appl.\ Mech.\ Engrg.\ 99 (1992) 209--233.

\bibitem{WarburtonHesthaven2003}
T.~Warburton and J.~S.~Hesthaven,
\emph{On the constants in $hp$-finite element trace inverse
inequalities},
Comput.\ Methods Appl.\ Mech.\ Engrg.\ 192 (2003) 2765--2773.

\bibitem{Bazilevs2007}
Y.~Bazilevs, V.~M.~Calo, J.~A.~Cottrell, T.~J.~R.~Hughes, A.~Reali and
G.~Scovazzi,
\emph{Variational multiscale residual-based turbulence modeling for
large eddy simulation of incompressible flows},
Comput.\ Methods Appl.\ Mech.\ Engrg.\ 197 (2007) 173--201.

\bibitem{Codina2000}
R.~Codina,
\emph{Stabilization of incompressibility and convection through
orthogonal sub-scales in finite element methods},
Comput.\ Methods Appl.\ Mech.\ Engrg.\ 190 (2000) 1579--1599.

\bibitem{Codina2002}
R.~Codina,
\emph{Stabilized finite element approximation of transient
incompressible flows using orthogonal subscales},
Comput.\ Methods Appl.\ Mech.\ Engrg.\ 191 (2002) 4295--4321.

\bibitem{CodinaVMSChapter}
R.~Codina, S.~Badia, J.~Baiges and J.~Principe,
\emph{Variational multiscale methods in computational fluid dynamics},
in: Encyclopedia of Computational Mechanics, 2nd ed., Wiley, 2018.

\end{thebibliography}

\end{document}
